% Optional standalone addition to Part E.  Sources have exact URLs in the
% accompanying note A_E_attempts.md, Section 12.  No shared bibliography edits.
\subsection{The norm and mixed-identity hypotheses do not suffice}

\begin{proposition}\label{prop:tree-norm-boundary}
For every countable group $K$ there is a countable nonabelian simple group
$E$ containing $K$ such that $E$ is not a torsion group, satisfies a
nontrivial mixed identity, and has
\[
 N_E(f)=\sup_{g\in E\setminus\{1\}}\rho_g^E(f)<\infty
 \quad(f\in E),\qquad \sup_{f\in E}N_E(f)=\infty.
\]
Moreover $E$ can be chosen to contain involutions $a,b$ for which the
unimodular equation
\begin{equation}\label{eq:tree-unsolvable-unimodular}
 xax a^{-1}x^{-1}=b
\end{equation}
has no solution.  In particular, taking $K$ to be Higman's group shows
that the other stated residual conditions do not imply uniform
simplicity when unimodular closure is omitted.
\end{proposition}

\begin{proof}
Consider the countable, normally colored tree $T$ with color code
\[
 \begin{pmatrix}
 0&2&3\\
 2&0&4\\
 \aleph_0&\aleph_0&0
 \end{pmatrix}.
\]
The entry in position $(i,j)$ is the number of color-$j$ neighbors
of a color-$i$ vertex.  Such a tree exists by the colored-tree code
criterion: zero entries are symmetric and the color-adjacency graph
is connected.  Normality means transitivity of the full
color-preserving automorphism group $A$ on each color class.  This
criterion and the following applications of Tits' theorem are recorded
in Gismatullin, \emph{Boundedly simple groups of automorphisms of trees},
Section~2.\footnote{\url{https://arxiv.org/pdf/0905.0913}, version~6,
Proposition~2.8 and the code discussion on pp.~6--8.  Journal DOI:
\url{https://doi.org/10.1016/j.jalgebra.2013.06.023}.}

Let $G=A^+$ be generated by the pointwise edge stabilizers.  All
nonzero code entries are at least two.  Thus $A$ has Tits'
independence property, preserves no proper nonempty subtree or end,
and $G$ is a nontrivial simple normal subgroup with these same
minimality properties.  Its action is of general type.

Write $\ell(f)$ for the least number of factors from pointwise
half-tree stabilizers needed to express $f\in G$.  This is finite,
because an edge stabilizer splits as a product of the two
complementary half-tree stabilizers.  Gismatullin's Proposition~2.8
gives, for every $g\ne1$, a bound of four on the $\rho_g^G$-length
of every element of every half-tree stabilizer.  If $g$ is
hyperbolic the bound is two.  Otherwise a product of two conjugates
of $g$ is hyperbolic, giving the bound four.  Consequently
\begin{equation}\label{eq:tree-pointwise-bound}
 N_G(f)\leq4\ell(f)<\infty.
\end{equation}
On the other hand, Theorem~3.11 of that paper would force $T$ to be
a uniform subdivision of a biregular tree if $G$ were uniformly
simple.  Its three vertex degrees $5,6,\aleph_0$, with no degree-two
vertices, rule this out.  Choose $f_m,g_m\in G$, $g_m\ne1$, with
$\rho_{g_m}^G(f_m)>m$ for every $m\geq1$.

At a color-two vertex there are countably infinitely many mutually
isomorphic rooted branches with a color-zero first vertex.  Fix
coherent identifications of these branches.  Permuting them
arbitrarily and fixing all other branches gives a copy of
$\operatorname{Sym}(\mathbb N)$ inside a pointwise edge stabilizer.
In particular $G$ contains a copy of the prescribed countable group
$K$.  Also choose coherent swaps $a=(12)$ and $b=(13)$ of three of
these branches, two transverse hyperbolic elements of $G$, and a
countable dense subset $D\subset G$ for the topology of pointwise
convergence on vertices.  Such $D$ exists because $T$ is countable
and $G$ is second countable.

Let $E_0$ be generated by all these countably many elements and all
the pairs $(f_m,g_m)$.  Given $E_n$, for every pair $f,g\in E_n$,
$g\ne1$, use \eqref{eq:tree-pointwise-bound} to choose an expression
of $f$ using at most $4\ell(f)$ conjugates of $g^{\pm1}$.  Adjoin
its conjugators for every such pair, obtaining the countable group
$E_{n+1}$.  Set $E=\bigcup_nE_n$.  Every pair in $E$ occurs at a
finite stage, so
\[
 \rho_g^E(f)\leq4\ell(f),\qquad N_E(f)<\infty.
\]
Thus every nonidentity element normally generates $E$, proving
simplicity.  The seeded pairs satisfy
$\rho_{g_m}^E(f_m)\geq\rho_{g_m}^G(f_m)>m$, proving unboundedness.
The seeded hyperbolic elements prove that $E$ is not torsion.

The faithful $E$-action on $T$ is minimal: density gives the same
vertex orbits as $G$, hence every $E$-invariant subtree is
$G$-invariant.  The seeded hyperbolic elements make the action of
general type, while $a$ pointwise fixes a half-tree.  Theorem~B of
Fima--Le Ma\^itre--Moon--Stalder therefore gives a nontrivial mixed
identity: for a faithful minimal general-type action of a countable
group on a tree, mixed-identity freeness is equivalent to topological
freeness on the boundary.\footnote{\emph{High transitivity for more
groups acting on trees}, Theorem~B, p.~4:
\url{https://webusers.imj-prg.fr/~pierre.fima/Documents/HTGroupsActingOnTrees.pdf}.}

It remains to verify the explicit failure of unimodular closure.
Put $c=ba$.  This cycles the three chosen branches and is the identity
on the complement.  Its fixed subtree has a unique vertex adjacent
to a nonfixed vertex, namely their common root $v$.  If $z^3=c$ in
$\operatorname{Aut}(T)$, then $z$ commutes with $c$ and preserves its
fixed subtree, so fixes $v$.  Its induced permutation $\sigma$ on
the neighbors of $v$ would have cube equal to a single three-cycle.
But $\sigma$ commutes with its cube, hence preserves the
three-element support of this cycle, and no permutation of three
elements cubes to a three-cycle.  Thus $c$ has no cube root even in
$\operatorname{Aut}(T)$.

The left side of \eqref{eq:tree-unsolvable-unimodular} equals
$(xa)x(xa)^{-1}$.  A solution would therefore be conjugate to the
involution $b$ and satisfy $x^2=1$.  Together with $a^2=1$, the
equation implies $(xa)^3=ba=c$, a contradiction.  Its exponent sum
in $x$ is one, as required.
\end{proof}

\paragraph{Scope.}
The last obstruction applies to subgroups containing the specified
coherent swaps.  It does not exclude arbitrary subgroups of these
tree automorphism groups.  In particular the proposition neither
constructs nor rules out a group satisfying the full residual
conjunction including unimodular closure.  The locally finite variant
with code
$\left(\begin{smallmatrix}0&2&3\\2&0&4\\3&4&0\end{smallmatrix}\right)$
already proves the same norm and mixed-identity boundary statement,
but its Higman images are trivial by the separate locally finite
tree obstruction.  Infinite valence is what permits the prescribed
subgroup $K$ in the version above.
